% !PRINTED 607
\weberpage{607}

% [corrected 2026-08-07, 10x zoom strip 1/10] The centred block here was the RUNNING HEAD of
% page 607 -- "§. 179. | Reelle Radicale. | 607" -- transcribed as a section heading. Running
% heads are excluded page furniture. The real §. 179 heading is on printed 606 (i0631) and reads
% "Auflösung durch reelle Radicale."
%
% Two tells, both usable without the scan: the number duplicates an existing heading, and the
% title is the SHORTENED form this book uses in running heads, not the full section title.
% Found by the section-series check; the second instance of this defect in Band I.
\quad Wir setzen einen \gesperrt{reellen Rationalitätsbereich} $\Omega$ voraus,
und nehmen in ihm eine Normalgleichung $g(t) = 0$ an,
d.\,h. eine irreducible Gleichung, deren Wurzeln alle rational
durch eine beliebige von ihnen, $\varrho$, ausdrückbar sind. Wenn eine
von den Wurzeln reell ist, so müssen auch alle anderen reell
sein, und $g(t)$ hat also entweder lauter reelle oder lauter paarweise
conjugirte imaginäre Wurzeln. $g(t) = 0$ kann die \gesperrt{Galois}'sche
Resolvente irgend einer gegebenen, sei es irreducibeln, sei es
reduciblen Gleichung $F(x) = 0$ sein, und wenn also $g(t)$ nur
reelle Wurzeln hat, so kann auch $F(x)$ nur relle Wurzeln haben, % [sic]
weil die Wurzeln von $F(x)$ rational durch eine Wurzel von $g(t)$
darstellbar sind. Wenn unter den Wurzeln von $F(x)$ imaginäre
sind, so hat $g(t)$ nur imaginäre Wurzeln. Hat aber $F(x)$
lauter reelle Wurzeln, so sind auch die Wurzeln von $g(t)$ reell,
weil sie ja rational durch die Wurzeln von $F(x)$ ausgedrückt
werden können.

\quad Wenn nun $g(t)$ durch Adjunction einer Wurzel $\varepsilon$ einer
irreducibeln Gleichung $\chi = 0$, \gesperrt{deren Grad eine Primzahl ist},
reducirt wird, so sind, wie wir im \S.\,157 gesehen haben, die
sämmtlichen Wurzeln von $\chi$ in $\Omega(\varrho)$ enthalten, und wenn also
$\varrho$ reell ist, so sind alle Wurzeln von $\chi$ reell.

\quad 1. \gesperrt{Eine Normalgleichung mit reellen Wurzeln kann
also nur durch solche irreducible Gleichungen
von Primzahlgrad reducirt werden, die lauter
reelle Wurzeln haben.}

\quad Wir fragen nun, ob eine Reduction der Normalgleichung
durch Adjunction eines reellen Radicales $\sqrt[p]{a}$ bewirkt werden
kann. Wir dürfen dabei annehmen, dass der Grad $p$ des Radicales
eine Primzahl sei, weil jedes Radical sich auf eine Reihe von
Wurzelziehungen von Primzahlgrad reduciren lässt. Auch können
wir voraussetzen, dass $a$ nicht die $p^{\text{te}}$ Potenz einer rationalen
Grösse sei, weil sonst die reelle Wurzel $\sqrt[p]{a}$ rational wäre.

\quad Unter dieser Vorraussetzung ist, wie wir jetzt beweisen wollen, % [sic]
$x^p - a$ irreducibel. Denn bezeichnen wir einen, z.\,B. den reellen
Werth $\sqrt[p]{a}$ mit $\alpha$ und eine imaginäre $p^{\text{te}}$ Einheitswurzel mit $\varepsilon$,
so sind
\weq{1}{\alpha,\ \varepsilon\alpha,\ \varepsilon^2\alpha \ldots \varepsilon^{p-1}\alpha}
%CONT>>
